\scanpage{123}
\mspage{7}
% archive sheet number: 122
\uline{NB} Il y a une notion duale de \uline{catégorie à carrés $h$-cocartésiens}, ou \uline{catégorie quadrangulée à droite}. Le plus souvent, les deux structures sont présentes simultanément. Il y a lieu de dégager des propriétés qui relient les deux structures duales sur une même catégorie. \inserted{Ainsi} Sauf erreur, dans le cas \inserted{$\unsure{\mathscr H|\mathcal H|H}$} \unsure{additife|additive|additi\deleted{f}ve|additif}, les deux espèces de carrés, les $h$-cartésiens et les $h$-cocartésiens, devraient \uline{coïncider}. (\unsure{Il|Je} me faudra revenir là-dessus par la suite.) Dans le cas $\mathscr H$ ``\unsure{ponctuée|pointée}'', il y a l'adjonction des foncteurs $\Sigma,\Omega$ liés l'un à la structure quadrangulée droite, l'autre à la structure quadrangulée gauche. Et il y a, quand il y a des $\underline{\operatorname{Hom}}$ internes ou externes, des propriétés éventuelles de $\underline{\operatorname{Hom}}(X,Y)$ \unsure{p.r.|p.v.} aux quadrangles $h$-cocartésiens en $X$, $h$-cartésiens en $Y$ -- mais c'est à peine un lien entre les deux structures --

2. Exemples.

Un premier exemple trivial, est celui où $\mathscr H$ est stable par produits fibrés, et où les carrés $h$-cartésiens sont définis comme les carrés cartésiens. Les axiomes $\operatorname{Cart}(1)$ à $\operatorname{Cart}(5)$ sont des propriétés bien connues, $\pm$ tautologiques de la notion de carré cartésien. Et
